% ============================================================
%   Abstract and Keywords  (English)
% ============================================================
%
% The class macros \hacerresumen and \resumenPersonalizado hard-code
% the Spanish word ``Resumen''. Because this document is written in
% English, the abstract is laid out manually here while the metadata
% macros \resumen and \descriptores (defined in main.tex) are still
% picked up by hyperref for PDF keywords.
% ------------------------------------------------------------

\cleardoublepage
\pagenumbering{roman}
\pagestyle{headings}

\vspace*{3cm}

{\bfseries\Large Abstract\par}
\vspace{0.5cm}

\noindent
This work presents CALA-SafeRL, a modular framework for Safe Reinforcement Learning applied to
autonomous driving in the CARLA simulator. The central challenge addressed is how to train a
Proximal Policy Optimization (PPO) agent that learns to drive under realistic traffic conditions
without catastrophic failures during the learning process itself. The framework is structured as a
composable chain of Gymnasium wrappers, allowing the agent, two interchangeable safety shields, a
dense reward shaper and the simulator to be combined with minimal coupling. The first shield
performs a lightweight threshold-based verification on semantic LIDAR readings together with
CARLA's Waypoint API; the second implements an adaptive-horizon strategy, predicting the vehicle's
future trajectory with a kinematic bicycle model calibrated for a Tesla Model 3 and selecting the
least invasive corrective action from a prioritised candidate set. The agent's observation space
has 739 dimensions, combining three semantic LIDAR channels with lane, lane-marking, vehicle and
route features derived from the Waypoint API. Training stability is improved through running-mean observation
normalisation, return scaling, Generalised Advantage Estimation, clipped surrogate updates with a
target-KL early stop, and a performance-driven curriculum that progressively introduces traffic
density. The framework is accompanied by a live metrics pipeline built on SQLite and a Streamlit
dashboard. Experiments are structured around the Town04 highway map and systematically compare the
three shielding configurations against baseline PPO. The resulting system is open, reproducible and
can serve as a foundation for further research into safety-aware policy learning.

\vspace{2cm}

\cleardoublepage
